% !TEX root = ../Projektdokumentation.tex
\section{Das Thema und Problemstellung}
\label{sec:ThemaUndProblemstellung}
Der Ausbildungsbetrieb betreibt und betreut Mail-Infrastrukturen für mehrere
Organisationen. In diesem Umfeld müssen freigegebene Mail-Adressen und
IP-Adressen regelmäßig angepasst werden, damit technische Ausnahmen im
Antispam-Umfeld nachvollziehbar und wirksam gepflegt bleiben. Vor Beginn des
Projekts erfolgte diese Pflege ohne zentrale Anwendung. Informationen waren auf
mehrere Datenquellen verteilt, Änderungen wurden manuell übernommen und
Rückfragen mussten häufig im direkten Austausch geklärt werden.

Die Ausgangssituation war fachlich problematisch, weil dieselben Daten an
verschiedenen Stellen gepflegt wurden und keine einheitliche Bedienlogik
vorhanden war. Dadurch entstanden Medienbrüche, uneinheitliche Datenstände
und ein zusätzlicher Abstimmungsaufwand zwischen den beteiligten
Mitarbeitern. Gleichzeitig war nur eingeschränkt nachvollziehbar, wer eine
Freigabe angelegt, geändert oder entfernt hatte. Für einen sicherheitsnahen
Fachprozess ist dieser Zustand langfristig nicht tragfähig.

\subsection{Zielstellung der Projektarbeit}
\label{sec:Zielstellung}
Ziel der Projektarbeit war die Entwicklung des
\Fachbegriff{AntispamAdminCenter}, einer webbasierten Verwaltungsanwendung für
die zentrale Pflege von Organisationen, Benutzern, Mail-Adressen und
IP-Adressen. Die Anwendung sollte von internen Administratoren und
autorisierten Organisationsbenutzern genutzt werden können und dabei eine
eindeutige Trennung zwischen globalen und mandantenbezogenen Rechten
sicherstellen.

Aus dem Projektauftrag ergaben sich folgende Kernziele:
\begin{itemize}
  \item zentrale und einheitliche Pflege aller fachlich relevanten Stammdaten
  \item rollenbasiertes Berechtigungskonzept mit Mandantentrennung
  \item abgesicherte Anmeldung mit Session-Verwaltung und Passwort-Hashing
  \item nachvollziehbare Protokollierung fachlicher Änderungen
  \item Bereitstellung einer belastbaren Datenbasis für nachgelagerte Mail-Systeme
\end{itemize}

Das Projektergebnis sollte nicht nur die Bearbeitungszeit senken, sondern vor
allem die Qualität der Datenhaltung verbessern. Neben der fachlichen Funktion
stand daher die Beherrschbarkeit des Betriebsprozesses im Vordergrund.

\subsection{Ist-Analyse}
\label{sec:IstAnalyse}
Die Ist-Analyse zeigte, dass die bisherige Vorgehensweise stark personenabhängig
war. Änderungen wurden zwar fachlich abgestimmt, technisch jedoch nicht über
eine dafür vorgesehene Anwendung umgesetzt. Daraus ergaben sich drei zentrale
Schwachstellen.

Erstens fehlte eine zentrale Datenhaltung für freigegebene Mail- und
IP-Einträge. Zweitens war die Berechtigungsstruktur nicht systematisch
abgebildet, so dass zwischen globalen Administratorrechten und
organisationsbezogenen Zuständigkeiten nur eingeschränkt unterschieden werden
konnte. Drittens existierte keine belastbare Historie, über die sich
Änderungen und deren Verursacher im Nachhinein nachvollziehen ließen.

Aus betrieblicher Sicht führte dies zu wiederkehrenden Rückfragen, manuellen
Kontrollen und vermeidbaren Fehlerquellen. Die eigentliche Pflegeaufgabe war
zwar fachlich überschaubar, der organisatorische Aufwand rund um Kontrolle,
Abstimmung und Nacharbeit fiel jedoch unverhältnismäßig hoch aus.

\subsection{Soll-Konzept}
\label{sec:SollKonzept}
Im Soll-Zustand soll eine zentrale Webanwendung den gesamten Pflegeprozess
unterstützen. Organisationen, Benutzer, Mail-Adressen und IP-Adressen werden
in einer gemeinsamen Datenbasis verwaltet und über einheitliche Formulare
gepflegt. Rollen und Organisationszuordnung steuern, welche Daten angezeigt und
geändert werden dürfen. Fachliche Änderungen werden revisionsfähig im
Audit-Log festgehalten.

Technisch wurde eine schlanke Webanwendung auf Basis von FastAPI, SQLModel und
SQLite vorgesehen. Diese Kombination war für den Projektumfang geeignet, weil
sie eine schnelle Umsetzung, klare Trennung zwischen Datenmodell, Fachlogik und
Webschicht sowie einfache Testbarkeit ermöglicht. Nicht Bestandteil des
Projekts sind die eigentliche Spam-Filterlogik, ein produktionsreifes
Mehrserver-Deployment oder die Migration beliebiger Altdatenbestände.

\subsection{Projektabgrenzung}
\label{sec:Projektabgrenzung}
Das Projekt umfasst die Entwicklung einer intern nutzbaren
Verwaltungsanwendung einschließlich Rollenmodell, Auditierung und
Kundendokumentation. Nicht Gegenstand der Projektarbeit sind der produktive
Betrieb einer hochverfügbaren Infrastruktur, eine allgemeine
Mandantenmigration aus Fremdsystemen oder die technische Umsetzung aller
denkbaren Schnittstellen zu nachgelagerten Mail-Systemen. Diese Abgrenzung war
erforderlich, um den Projektumfang auf einen in 80 Stunden realisierbaren und
bewertbaren Kern zu konzentrieren.
